\chapter{Particle movers, forces and field solvers}\label{ch:movers}
\section{Neutral and charged characteristics}
For a neutral under a prescribed force,
$\dot{\boldsymbol x}=\boldsymbol v$ and
$\dot{\boldsymbol v}=\boldsymbol F/m$.
Gravity from a spherical body is
$-\mathcal GM\,\boldsymbol r/r^3$ outside its mass distribution, with
$\mathcal GM$ in $\mathrm{m^3\,s^{-2}}$. Nonspherical gravity requires a
separately defined mass model; a surface shape alone does not determine it.
Radiation pressure and gas drag are distinct accelerations with their own
size, velocity and geometry dependence.

Charged-particle trajectories satisfy
\begin{equation}
\dot{\boldsymbol v}=\frac{q}{m}
(\boldsymbol E+\boldsymbol v\times\boldsymbol B)+\boldsymbol a_{\rm other}.
\end{equation}
The relativistic equation advances momentum
$\boldsymbol p=\gamma m\boldsymbol v$ with
$\gamma=\sqrt{1+p^2/(m^2c^2)}$; replacing $m$ by a constant guessed
relativistic mass is generally incorrect. Validate gyrofrequency, phase,
parallel acceleration and drift against independent solutions.

\section{Selecting a mover}
The ordinary host dispatch is
\code{PIC::Mover::DispatchParticleMover(ptr, dt, node)}.
Legacy macro selection remains the default. Pointer selection requires its
compiled legacy-settings switch to be off and a valid
\code{PIC::Mover::SetUserDefinedParticleMover} callback installed before
advancing a particle. Boundary-specific or device bulk kernels can have
separate selection paths. The callback owns deletion/list insertion/migration
according to the configured mover contract.

The timestep must resolve the fastest relevant gyroperiod and field scale,
as well as cell crossings. Guiding-center motion assumes gyroaveraging is
valid: gyroradius and gyroperiod must be small relative to background scales.
Its drifts, mirror force and parallel electric force are not a valid fallback
at a magnetic null; the field-local API provides typed invalid-domain status.

\section{Electromagnetic PIC and ECSIM}
Self-consistent PIC deposits particle charge/current and solves fields rather
than reading a prescribed background. ECSIM is an energy-conserving
semi-implicit formulation in the supported configuration. Verify the actual
spatial operators, mass/moment assembly, boundary rows, solver residuals and
particle-field energy accounting. Stability of a semi-implicit method does
not remove resolution requirements for the phenomenon of interest.

For deposited shape weights $S_i$, require the intended partition of unity
and $\sum_i Q_i=q w$ for one particle. Current conservation additionally
requires a compatible discrete continuity equation, not just a correct total
charge. Manufactured field tests establish consistency of operators on the
tested meshes; they are not a blanket validation of every production plasma.

\section{Boundary events and complete trajectories}
Find the earliest physical boundary intersection within the remaining
trajectory interval. Reflect or absorb at that event, then process the
remaining time exactly once. A particle reaching a remote block must be
delivered to its owner even if its continuation marker is already zero.
Repeated reflection, tangency and starts on/inside a boundary need explicit
tolerance and ownership conventions.

Use independent straight-line, constant-acceleration, uniform-field gyro,
moving-plane and moving-sphere solutions as test oracles. Check phase-space
invariants, event time, list membership, endpoint and represented losses;
checking only final radius misses wrong velocity and duplicate ownership.
\source{src/pic/README.md; test/srcCoreUnitTests/README.md; src/test/tests/pic/ecsim_operators/README.md}
